% ===== §4 Epistemology I: the quantum-structured grammar =====
\section{The Quantum-Structural Articulation of the Relational Grammar}
\label{sec:grammar}

\noindent
The ontology of \xref{sec:ontology} establishes four conditions: value is not antecedent to its co-apprehension; apprehension is identically generation; the manner of apprehension conditions its term; and the dynamics is rewritten by the term it generates. A cognition satisfying these conditions cannot be representational, since representation presupposes an antecedent term; it is participatory, in the sense that the cognition is a component of the generation of what is cognised. The present section establishes that the formal structure of quantum measurement supplies a precise articulation of participatory cognition, and writes the relational grammar in that structure. The articulation is structural and epistemic, and its scope is delimited at the outset.

\subsection{Delimitation of the claim}
\label{sec:self-limit}

The structure of quantum theory is employed here as a formal model of participatory generation. No claim is advanced that intimate relation, joint attention, or the neural substrate is a quantum-physical system; no hypothesis is made concerning quantum coherence, entanglement, or physical collapse in any relational or cognitive substrate. The positions that do advance such claims, orchestrated objective reduction, and the literalist interpretation of quantum cognition, are not at issue, and the present argument is independent of their physical adjudication. The procedure is continuous with the series' prior use of probabilistic context-free grammars to model the recursion of generative justice \citep{eglash2016,eglash1999}: the assertion that a regenerative structure has the form of a branching process is an assertion about form and not about physical constitution, and the assertion that the co-apprehension of value has the form of a quantum measurement is of the same kind.

\begin{claim}[The quantum structure as epistemic model]
Quantum structure is employed to articulate, in formal terms, the participatory character of co-generation established phenomenologically in \xref{sec:structure} and required ontologically in \xref{sec:ontology}. The relation asserted is a structural isomorphism at the level of epistemology and not a thesis concerning physical substrate. No claim of the paper depends upon relation or cognition being physically quantum; the quantum vocabulary is, in principle, replaceable by ``the formal structure of participatory generation,'' for which quantum measurement is the developed exemplar. At the single point at which the structural model and the physics diverge, the question of enumerability, \xref{sec:unenumerable}, the structure is retained and the physical thesis explicitly declined.
\end{claim}

\subsection{The relational state as superposition}
\label{sec:state}

Let the value-potential of a relation at a given moment be represented by a state $|\psi\rangle$ in a complex inner-product space $\mathcal{H}$, the relational state space. Let $\{|v_i\rangle\}$ be an orthonormal family of \emph{determinate value-states}, the values the relation might come to bear, each $|v_i\rangle$ a candidate determinate value. The relational state is in general a superposition,
\[
|\psi\rangle \;=\; \sum_i c_i\,|v_i\rangle, \qquad \sum_i |c_i|^2 = 1,
\]
in which no single $|v_i\rangle$ is the value the relation determinately bears. This represents formally the ontological thesis that real value is not antecedent to its co-apprehension (\xref{sec:opening}). The interpretation of $|\psi\rangle$ is to be distinguished from a probability distribution over antecedently determinate values: the coefficients $c_i$ do not encode ignorance of a value already settled but the genuine indeterminacy of a value not yet generated, for which no further specification of the present state resolves the indeterminacy, there being as yet no determinate term. The opening of \xref{sec:opening}, the site at which something is not yet determinate, is formally the superposition prior to its resolution.

\subsection{Co-apprehension as participatory measurement}
\label{sec:measurement}

Joint attention is modelled not as a registration of $|\psi\rangle$ but as a \emph{measurement} of it. A measurement is specified by a choice of observable, equivalently by a choice of orthonormal basis $\{|b_j\rangle\}$ of $\mathcal{H}$ in which the measurement is resolved; the basis is not fixed by the state but is selected, and the selection corresponds to the manner of apprehension. This is the formal content of the thesis that apprehension is generation (\xref{sec:hinge}): the measurement does not return an antecedent value but resolves the superposition into a determinate term. The decisive structural fact is the basis-dependence of the outcome. Expressed in the basis $\{|b_j\rangle\}$, the state is $|\psi\rangle = \sum_j d_j |b_j\rangle$ with $d_j = \langle b_j | \psi\rangle$, and the values accessible as outcomes are exactly the $|b_j\rangle$; a distinct choice of basis $\{|b'_k\rangle\}$ renders accessible a distinct family of outcomes, and in general $[\,\hat{B}, \hat{B}'\,] \neq 0$, so that the two families are not jointly resolvable.

\begin{claim}[The manner of apprehension as choice of basis]
The manner in which two parties jointly attend corresponds formally to the selection of a measurement basis, and the basis determines the family of values that can be co-generated. There is no basis-independent measurement, and hence no manner of joint attention that merely registers an antecedent value; every manner of attending is the selection of a space of admissible outcomes, effected prior to the determination of any outcome. The phenomenological thesis that the manner conditions the term (\xref{sec:hinge}) is thereby given formal content: the manner is the basis, and the basis fixes the space of admissible terms.
\end{claim}

\subsection{Generation as collapse with reconfiguration of the state space}
\label{sec:collapse}

The generation of a determinate value is modelled by two operations conjoined. The first is resolution: the superposition $|\psi\rangle$ is resolved, on this occasion, into a determinate outcome $|b_j\rangle$. The second is reconfiguration: the post-measurement relational state space is not identical to the prior space. Where standard measurement returns the system to a state $|b_j\rangle$ within a fixed $\mathcal{H}$, the present model requires that the admissible state space itself be altered by the resolution, the realised value $|b_j\rangle$ enters as a condition on subsequent dynamics, so that the family of admissible subsequent superpositions, and in the limiting case the dimension of the effective state space, is transformed. Formally, the realised value induces a map $\mathcal{H} \to \mathcal{H}'$ in which dimensions are added (relational possibilities that did not antecedently exist) or removed (possibilities foreclosed). This is the formal content of the reflexive thesis of \xref{sec:reflexive}: the generated term reconfigures the dynamics that bore it.

\begin{claim}[Generation as collapse-cum-reconfiguration]
Each co-generation of value is conjointly a resolution of the relational superposition into a determinate term and a transformation of the space of subsequently admissible terms. The generated value is not the terminus of a fixed dynamics but a condition that reconfigures the dynamics, altering the admissible superpositions available for subsequent generation. Measurement modifies the measured system; in the relational case the modification is constitutive of generation rather than incidental to it.
\end{claim}

\subsection{Geometric phase}
\label{sec:phase-source}

The geometric phase employed throughout the series to characterise the good cycle, the return to a configuration accompanied by a non-trivial accumulated holonomy, distinguished from mere repetition, has its origin in the quantum structure here invoked. A state transported around a closed loop in the projective state space accumulates a phase determined by the geometry of the loop and not by its duration \citep{berry1984,pancharatnam1956,aharonov1987}. For a loop $\gamma$ in the parameter space governing the relational state, the accumulated geometric phase is
\[
\phi_{\mathrm{geo}}(\gamma) \;=\; \oint_\gamma \langle \psi \,|\, i\,\nabla \,|\, \psi \rangle \cdot d\mathbf{R},
\]
the holonomy of the natural connection on the state bundle. The anchoring of the relational grammar in quantum structure thus returns the series' geometric phase to its source rather than importing a foreign formalism. The model further specifies a feature the prior treatment could only assert: since the measurement basis is variable (\xref{sec:measurement}), the connection along which phase is accumulated is itself variable, and a variation in the manner of joint attention is formally a variation of the connection, hence of the holonomy that a given loop in relational configuration accumulates. Two relations may traverse the same configurations and accumulate distinct, even opposite, phases, in virtue of attending, connecting, differently. The Value-Foam programme is, on this articulation, the determination of the carrier and the connection of which the series' relational holonomy is an instance, with Berry's phase as its formal antecedent.

\subsection{Enumerability and non-possession}
\label{sec:unenumerable}

The ontology established that a system whose dynamics is rewritten by its products admits no external pre-enumeration (\xref{sec:reflexive}). The quantum articulation renders this precise: if each co-generation conjointly resolves a superposition and reconfigures the space of subsequent superpositions, then no grammar fixed in advance, external to the relation, enumerates the values the relation will generate, since the productions governing subsequent generation are themselves generated within the process and are not antecedently available for enumeration.

The strength of this claim requires delimitation, in accordance with \xref{sec:self-limit}. The claim is not that the relational grammar exceeds the recursively enumerable in the computability-theoretic sense, that its generation computes a non-recursive function. That thesis is not advanced and is not supported by the structure: quantum computation does not compute non-recursive functions, and quantum indeterminacy yields unpredictability rather than super-recursive enumerative capacity. The non-enumerability asserted is epistemic and internal.

\begin{claim}[Non-enumerability as the form of non-possession]
The relational grammar is non-enumerable in the following sense: there exists no grammar, fixed in advance and external to the relation, that enumerates the values the relation will co-generate, since the generating productions are produced within the joint act that any external enumeration would be required to anticipate. The complete fixing of the grammar is not a tractable computation that happens to be infeasible; it is the substitution of a settled rule-set for a generation in progress, which is the operation of settling and of possession. Enumerability, settledness, and possession are coextensive, and their negation is the formal correlate of the non-possessibility of the relation. The impossibility of fixing the grammar is therefore not an epistemic limitation but the formal signature of the openness the phenomenology described as the sustaining of the nameless and the ontology declined to complete.
\end{claim}

\noindent
The quantum articulation thus discharges, in formal terms, the conditions established by the phenomenology and the ontology, and no more: a state of superposed potential in place of a determinate value; a co-apprehension modelled as participatory measurement with a basis selected in the manner of attending; a generation modelled as collapse conjoined with reconfiguration of the state space; a geometric phase returned to its quantum source, with the variable basis as the formal correlate of the efficacy of manner; and a non-enumerability that is the epistemic correlate of non-possession, with the computability-theoretic thesis declined. The articulation constitutes one realisation of the reflexive structure, its formal-epistemic phase. The following section develops the same structure under the conditions of political economy, where it is realised in a distinct form: the circulation of value that reconfigures its subject.
